\documentclass[a4paper]{article}

% Basic packages
% Español (México) con XeLaTeX: fontspec para fuentes Unicode, polyglossia
% para la división silábica y los nombres del documento en la variante mexicana.
\usepackage{fontspec}
\usepackage{polyglossia}
\setdefaultlanguage[variant=mexican]{spanish}
% Los nombres chinos van como «pinyin (original)»: xeCJK da a los caracteres
% Han su propia fuente; sin ella XeLaTeX los omite con un aviso.
\usepackage{xeCJK}
\setCJKmainfont{FandolSong-Regular.otf}
% polyglossia no traduce los nombres de listings.
\AtBeginDocument{\providecommand{\lstlistingname}{}\renewcommand{\lstlistingname}{Listado}%
  \providecommand{\lstlistlistingname}{}\renewcommand{\lstlistlistingname}{Índice de listados}}
\usepackage{amsmath, amssymb}
\usepackage{graphicx}
\usepackage[margin=2.5cm]{geometry}
\usepackage[most]{tcolorbox}
\usepackage{etoolbox}
\usepackage{listings}
\usepackage{hyperref}
\usepackage{booktabs}
\usepackage{subcaption}
\usepackage{float}

% Highlight boxes
\newtcolorbox{knowledgebox}[1]{
    enhanced, colback=blue!5!white, colframe=blue!75!black, colbacktitle=blue!75!black,
    coltitle=white, fonttitle=\bfseries, title={#1}, attach boxed title to top left={yshift=-2mm, xshift=2mm},
    boxrule=1pt, sharp corners
}

\newtcolorbox{importantbox}[1]{
    enhanced, colback=yellow!10!white, colframe=yellow!80!black, colbacktitle=yellow!80!black,
    coltitle=black, fonttitle=\bfseries, title={#1}, sharp corners
}

\newtcolorbox{warningbox}[1]{
    enhanced, colback=red!5!white, colframe=red!75!black, colbacktitle=red!75!black,
    coltitle=white, fonttitle=\bfseries, title={#1}, sharp corners
}

% Listings
\lstset{
    language=Python,
    basicstyle=\ttfamily\small,
    keywordstyle=\color{blue},
    stringstyle=\color{red!60!black},
    commentstyle=\color{green!60!black},
    breaklines=true,
    frame=single,
    numbers=left,
    numberstyle=\tiny\color{gray},
    captionpos=b,
    extendedchars=true
}

% Front-page metadata
\newcommand{\notetitle}{LightRAG avanzado: KG-based RAG y visualización de grafos de conocimiento}
\newcommand{\noteauthors}{Elaboradas a partir de materiales públicos del curso}
\newcommand{\notedate}{2025}
\newcommand{\videochannel}{Wudaokou Naishi (五道口纳什)}
\newcommand{\videopublishdate}{2025}
\newcommand{\videoduration}{aproximadamente 12 minutos}
\newcommand{\videourl}{}
\newcommand{\videocoverpath}{cover.jpg}
\newcommand{\repourl}{https://github.com/hqhq1025/ai-course-notes}


% Footer with repo info
\usepackage{fancyhdr}
\pagestyle{fancy}
\fancyhf{}
\fancyfoot[C]{\thepage}
\fancyfoot[R]{\footnotesize\color{gray}\href{\repourl}{AI Course Notes}}
\renewcommand{\headrulewidth}{0pt}
\renewcommand{\footrulewidth}{0pt}

\begin{document}

\begin{titlepage}
\centering
{\Large Notas del curso\par}
\vspace{1.2cm}
{\huge\bfseries \notetitle\par}
\vspace{0.8cm}
{\large \noteauthors\par}
\vspace{0.3cm}
{\large \notedate\par}
\vspace{1.2cm}

\ifdefempty{\videocoverpath}{
{\small Al generar las notas, indique la ruta local de la portada del video.\par}
}{
\includegraphics[width=0.82\textwidth,height=0.45\textheight,keepaspectratio]{\videocoverpath}\par
}

\vfill
\begin{tcolorbox}[width=0.9\textwidth, colback=black!2!white, colframe=black!60, sharp corners]
\textbf{Autor o canal del video}: \videochannel\par
\textbf{Fecha de publicación}: \videopublishdate\par
\textbf{Duración del video}: \videoduration\par
\ifdefempty{\videourl}{}{
\textbf{Enlace del video}: \href{\videourl}{\nolinkurl{\videourl}}\par
}
\par
\textbf{Página del proyecto}: \href{\repourl}{\nolinkurl{\repourl}}\par
\end{tcolorbox}
\end{titlepage}

\tableofcontents
\newpage

%% --- Inicio del contenido --- %%

\section{Introducción}

Esta sesión es la continuación avanzada de LightRAG, que retoma el contenido de la sesión anterior y profundiza en los detalles del código fuente y el mecanismo de funcionamiento de LightRAG. Cubre principalmente cuatro aspectos:

\begin{enumerate}
    \item Por qué introducir un RAG basado en Knowledge Graph Base
    \item El proceso de indexación incremental (Mode) de LightRAG
    \item La visualización del Knowledge Graph
    \item Una explicación detallada de los modos de recuperación Hybrid/Mix
\end{enumerate}


\section{Por qué se necesita KG-based RAG}

\subsection{Limitaciones del Naive RAG}

El RAG tradicional basado en Chunks (Naive) tiene dos limitaciones centrales:

\begin{importantbox}{Los dos puntos débiles del Naive RAG}
\begin{itemize}
    \item \textbf{Dificultad con el razonamiento multisalto}: cuando la relación entre entidades cruza varios Chunks (A→B→C→D→E), el Naive RAG solo puede recuperar fragmentos locales y no logra enlazar la cadena de razonamiento completa
    \item \textbf{Ausencia de búsqueda global}: al estar limitado por la recuperación TopK, si el Chunk clave no está entre los TopK, el razonamiento falla
\end{itemize}
\end{importantbox}

\subsection{Ventajas del KG-based RAG}

El RAG basado en Knowledge Graph difiere del Naive RAG de manera esencial, tanto en la forma de los datos como en la manera de razonar:

\begin{table}[H]
\centering
\begin{tabular}{p{3cm}p{5cm}p{5cm}}
\toprule
\textbf{Dimensión} & \textbf{Naive RAG (basado en Chunks)} & \textbf{KG-based RAG} \\
\midrule
Forma de los datos & Puntos (un Embedding por cada Chunk) & Red (Graph formado por entidades + relaciones) \\
Riqueza semántica & Texto original & Las entidades/relaciones tienen descripciones detalladas y Embedding \\
Manera de razonar & Recuperar Chunks → colocarlos en el Context → razonamiento fragmentado & Recorrido de la estructura del grafo + descubrimiento de asociaciones → agregación estructurada \\
Alcance global & Limitado por TopK & Búsqueda global mediante el recorrido del grafo \\
Interpretabilidad & Baja & Alta (la ruta se puede rastrear) \\
\bottomrule
\end{tabular}
\caption{Comparación central entre Naive RAG y KG-based RAG}
\end{table}

\begin{knowledgebox}{La construcción del KG depende por completo del modelo de lenguaje}
En LightRAG, la extracción de entidades/relaciones y la generación de descripciones (Profiling) se realizan enteramente mediante Prompts al modelo de lenguaje. Esto significa que la calidad del KG depende directamente de la capacidad del LLM.
\end{knowledgebox}

\subsection{Resumen de la sección}

El KG-based RAG transforma los documentos en un grafo de conocimiento estructurado, lo que permite el razonamiento multisalto y la búsqueda global, y compensa las carencias del RAG tradicional basado en Chunks en escenarios de razonamiento complejo.

\section{Proceso de indexación incremental (Modo)}

\subsection{Flujo de procesamiento de documentos nuevos}

LightRAG admite la incorporación dinámica de documentos nuevos; su proceso de Modo (Upsert = Update + Insert) es el siguiente:

\begin{enumerate}
    \item \textbf{Extracción de entidades y relaciones}: con base en un modelo de lenguaje, se extraen entidades (Entity) y relaciones (Edge) del documento nuevo
    \item \textbf{Profiling}: se genera una descripción textual detallada para cada entidad y relación
    \item \textbf{Embedding}: se construye un índice vectorial para las descripciones de entidades, las descripciones de relaciones y el texto de los Chunks
    \item \textbf{Fusión de subgrafos}: el subgrafo del documento nuevo se fusiona con el grafo de conocimiento global
\end{enumerate}

\subsection{Alineación de entidades y consistencia de nombres}

\begin{warningbox}{La consistencia en el nombramiento de entidades es clave}
Una misma entidad puede mencionarse de formas distintas en diferentes Chunks. LightRAG lo especifica explícitamente en el System Prompt de Entity Extraction:
\begin{itemize}
    \item Los nombres de entidades \textbf{no distinguen mayúsculas de minúsculas}
    \item Cada palabra con significado debe escribirse con mayúscula inicial (por ejemplo, «Knowledge Graph» y no «knowledge graph»)
    \item Debe garantizarse un nombramiento consistente (la misma entidad usa siempre el mismo nombre)
\end{itemize}
La alineación de entidades depende principalmente de la consistencia del modelo de lenguaje al generar el texto, y no de un emparejamiento posterior.
\end{warningbox}

\subsection{Estrategia de fusión de descripciones: MapReduce}

Cuando una misma entidad o relación aparece en varios Chunks, cada Chunk aporta una \textbf{descripción parcial} de esa entidad o relación. La estrategia de fusión sigue el patrón MapReduce:

\begin{enumerate}
    \item \textbf{Recolección}: se reúnen todas las descripciones extraídas de los distintos Chunks
    \item \textbf{Evaluación de longitud}:
    \begin{itemize}
        \item Descripción corta $\rightarrow$ se concatena directamente con un separador
        \item Descripción larga $\rightarrow$ se comprime semánticamente (Summarization) con base en un modelo de lenguaje
    \end{itemize}
    \item \textbf{Actualización sincronizada}: \texttt{src\_id} también se concatena, para registrar el origen de la descripción
\end{enumerate}

\begin{importantbox}{Prompt de compresión semántica}
En \texttt{prompts.py} de LightRAG se define el Prompt \texttt{summarize\_entity\_descriptions}, dedicado a comprimir descripciones de entidades o relaciones demasiado extensas en un resumen más conciso, conservando la información esencial.
\end{importantbox}

\subsection{Ejemplo de fusión}

Tomemos dos Chunks como ejemplo:
\begin{itemize}
    \item \textbf{Chunk 1}: se extraen la entidad Alex (descripción A), Bob (descripción A) y la relación Alex→Bob = «colegas»
    \item \textbf{Chunk 2}: se extraen la entidad Alex (descripción B) y la relación Alex→Bob = «amigos»
\end{itemize}

Después de la fusión:
\begin{itemize}
    \item Descripción de Alex = descripción A \texttt{<sep>} descripción B (superposición de descripciones parciales)
    \item Descripción de la relación Alex→Bob = «colegas» \texttt{<sep>} «amigos»
    \item Si la longitud total de la descripción supera el umbral, se activa la compresión semántica
\end{itemize}

\subsection{Resumen de la sección}

La indexación incremental de LightRAG logra la incorporación fluida de documentos nuevos mediante tres pasos: alineación de entidades, fusión de descripciones (MapReduce) y compresión semántica. Todo el proceso depende en gran medida de la capacidad del modelo de lenguaje.
\section{Visualización del grafo de conocimiento}

LightRAG genera durante el proceso de indexación un archivo de grafo en formato GraphML. El proyecto ofrece código de visualización basado en la biblioteca \texttt{pyvis}:

\begin{itemize}
    \item Al hacer clic en un nodo se puede ver la descripción detallada de la entidad
    \item Las aristas muestran la información descriptiva de la relación
    \item Es adecuado para depurar y entender la estructura del KG
\end{itemize}

\begin{knowledgebox}{Formato GraphML}
GraphML es un formato de descripción de estructuras de grafo basado en XML, ampliamente usado para almacenar e intercambiar grafos de conocimiento. LightRAG genera este archivo automáticamente en la etapa de Index, sin necesidad de configuración adicional.
\end{knowledgebox}

\section{Modos de recuperación en detalle: Naive, Hybrid, Mixed}

\subsection{Los tres modos de recuperación}

LightRAG admite tres modos de recuperación (Query Mode), cada uno con un caso de uso propio:

\begin{importantbox}{Comparación de los tres modos de recuperación}
\begin{itemize}
    \item \textbf{Naive}: se basa únicamente en Vector — se calcula el Embedding de la Query, se obtiene la Cosine Similarity con el Embedding de cada Chunk y se devuelve el TopK
    \item \textbf{Hybrid}: se basa por completo en el grafo de conocimiento — combina Local Search (vecindad de entidades) y Global Search (coincidencia de relaciones), sin consultar directamente los fragmentos de texto originales en la base vectorial
    \item \textbf{Mixed}: Hybrid + Naive — combina el razonamiento global del KG con la recuperación vectorial tradicional para obtener los fragmentos de texto originales
\end{itemize}
\end{importantbox}

\subsection{Extracción de palabras clave: Low-Level y High-Level}

La recuperación de LightRAG primero extrae dos tipos de palabras clave de la Query del usuario:

\begin{itemize}
    \item \textbf{Palabras clave Low-Level}: nombres concretos de entidades (como nombres de personas o de productos), usadas para Local Search
    \item \textbf{Palabras clave High-Level}: temas macro, conceptos abstractos, usadas para Global Search
\end{itemize}

Este paso también lo realiza el modelo de lenguaje, y su entrada y salida se pueden monitorear con el Trace de LangFuse.

\subsection{Local Search (búsqueda local)}

\begin{enumerate}
    \item Se usan las palabras clave Low-Level para hacer Vector Search en la base vectorial de entidades
    \item Se recuperan los nodos de entidad correspondientes
    \item Se hace un recorrido (Traverse) de un salto en el Graph para obtener las entidades vecinas directamente relacionadas
    \item Se devuelven las entidades y relaciones dentro de ese salto como Context
\end{enumerate}

Local Search está centrada en las entidades y es adecuada para responder preguntas detalladas sobre una entidad específica.

\subsection{Global Search (búsqueda global)}

\begin{enumerate}
    \item Se usan las palabras clave High-Level para hacer Vector Search en la base vectorial de relaciones
    \item Se recuperan las relaciones (aristas) correspondientes
    \item Se devuelven esas relaciones junto con las entidades asociadas como Context
\end{enumerate}

Global Search est�� centrada en las relaciones, se enfoca en temas macro y es adecuada para responder preguntas integrales que abarcan varios dominios o documentos.

\subsection{Resumen de la sección}

LightRAG logra una recuperación multinivel, desde entidades concretas hasta temas macro, mediante la separación de palabras clave Low-Level/High-Level y la doble búsqueda Local/Global. El modo Mixed, además, incorpora los fragmentos de texto originales, con lo que ofrece la cobertura de información más completa.

\section{Síntesis y ampliación}

En esta entrega se analizaron a fondo cuatro detalles técnicos clave de LightRAG:

\begin{enumerate}
    \item \textbf{La necesidad de KG-based RAG}: resuelve las limitaciones fundamentales de Naive RAG en el razonamiento multisalto y las consultas globales
    \item \textbf{Mecanismo de indexación incremental}: mediante fusión de descripciones al estilo MapReduce y compresión semántica, permite agregar documentos de forma dinámica
    \item \textbf{Visualización del KG}: ofrece una herramienta intuitiva de exploración del grafo basada en pyvis
    \item \textbf{Recuperación multimodo}: los tres modos Naive/Hybrid/Mixed cubren distintas necesidades de consulta
\end{enumerate}

\subsection{Lecturas adicionales}

\begin{itemize}
    \item El archivo \texttt{prompts.py} del código fuente de LightRAG: contiene todos los prompts clave (Entity Extraction, Summarization, etc.)
    \item Despliegue local de LangFuse: para monitorear las entradas y salidas de todas las llamadas a la API de LightRAG
    \item Microsoft GraphRAG: una solución de KG-based RAG similar a LightRAG, útil para un estudio comparativo
\end{itemize}

%% --- Fin del contenido --- %%

\end{document}
